\documentclass[12pt,a4paper]{article}

% Encoding and fonts
\usepackage[utf8]{inputenc}
\usepackage[T1]{fontenc}
\usepackage{lmodern}
\usepackage{textcomp}

% Page layout
\usepackage[top=1in, bottom=1in, left=1in, right=1in]{geometry}

% Typography
\usepackage{microtype}
\usepackage{parskip}

% Language
\usepackage[english]{babel}

% Headers and footers
\usepackage{fancyhdr}
\pagestyle{fancy}
\fancyhf{}
\fancyhead[L]{\leftmark}
\fancyhead[R]{\thepage}
\fancyfoot[C]{\thepage}
\renewcommand{\headrulewidth}{0.4pt}
\renewcommand{\footrulewidth}{0pt}

% Section styling
\usepackage{titlesec}
\titleformat{\section}{\Large\bfseries}{\thesection}{1em}{}
\titleformat{\subsection}{\large\bfseries}{\thesubsection}{1em}{}
\titleformat{\subsubsection}{\normalsize\bfseries}{\thesubsubsection}{1em}{}
\titlespacing*{\section}{0pt}{2.5ex plus 1ex minus .2ex}{1.5ex plus .2ex}
\titlespacing*{\subsection}{0pt}{2ex plus 1ex minus .2ex}{1ex plus .2ex}
\titlespacing*{\subsubsection}{0pt}{1.5ex plus 1ex minus .2ex}{0.8ex plus .2ex}

% List formatting
\usepackage{enumitem}
\setlist[itemize]{topsep=0.5ex, itemsep=0.25ex, parsep=0pt}
\setlist[enumerate]{topsep=0.5ex, itemsep=0.25ex, parsep=0pt}

% Essential packages
\usepackage{graphicx}
\usepackage[table]{xcolor}
\usepackage{booktabs}
\usepackage{array}
\usepackage{tabularx}
\usepackage{longtable}
\usepackage{amsmath}
\usepackage{amssymb}
\usepackage[normalem]{ulem}
\rowcolors{2}{gray!10}{white}
\renewcommand{\arraystretch}{1.3}

% Cross-references
\usepackage{hyperref}
\usepackage{cleveref}

% Custom commands
\newcommand{\pilcrow}{\S}

% Hyperref setup
\hypersetup{
    colorlinks=true,
    linkcolor=blue,
    filecolor=magenta,
    urlcolor=cyan,
}

% Code listings
\usepackage{listings}
\definecolor{codebg}{RGB}{248,248,248}
\definecolor{codeframe}{RGB}{200,200,200}
\definecolor{codekeyword}{RGB}{0,0,180}
\definecolor{codecomment}{RGB}{0,128,0}
\definecolor{codestring}{RGB}{163,21,21}
\lstset{
    basicstyle=\ttfamily\small,
    breaklines=true,
    breakatwhitespace=false,
    frame=single,
    framerule=0.5pt,
    rulecolor=\color{codeframe},
    backgroundcolor=\color{codebg},
    keepspaces=true,
    numbers=left,
    numberstyle=\tiny\color{gray},
    numbersep=8pt,
    showstringspaces=false,
    tabsize=4,
    xleftmargin=1em,
    xrightmargin=1em,
    aboveskip=1em,
    belowskip=1em,
    keywordstyle=\color{codekeyword}\bfseries,
    commentstyle=\color{codecomment}\itshape,
    stringstyle=\color{codestring},
    literate={_}{{\textunderscore}}1
             {-}{{\textminus}}1
             {<}{{\textless}}1
             {>}{{\textgreater}}1
             {~}{{\textasciitilde}}1
             {^}{{\textasciicircum}}1
}

% YAML language definition for listings
\lstdefinelanguage{YAML}{
    keywords={true,false,null,y,n},
    keywordstyle=\color{blue}\bfseries,
    sensitive=false,
    comment=[l]{\#},
    morecomment=[s]{/*}{*/},
    commentstyle=\color{gray}\ttfamily,
    stringstyle=\color{red}\ttfamily,
    morestring=[b]'',
    morestring=[b]'',
}

% TOML language definition for listings
\lstdefinelanguage{TOML}{
    keywords={true,false},
    keywordstyle=\color{blue}\bfseries,
    sensitive=true,
    comment=[l]{\#},
    commentstyle=\color{gray}\ttfamily,
    stringstyle=\color{red}\ttfamily,
    morestring=[b]'',
    morestring=[b]'',
}

% Dockerfile language definition for listings
\lstdefinelanguage{Dockerfile}{
    keywords={FROM,RUN,CMD,LABEL,EXPOSE,ENV,ADD,COPY,ENTRYPOINT,VOLUME,USER,WORKDIR,ARG,ONBUILD,STOPSIGNAL,HEALTHCHECK,SHELL},
    keywordstyle=\color{blue}\bfseries,
    sensitive=false,
    comment=[l]{\#},
    commentstyle=\color{gray}\ttfamily,
    stringstyle=\color{red}\ttfamily,
    morestring=[b]'',
    morestring=[b]'',
}

% JSON language definition for listings
\lstdefinelanguage{JSON}{
    keywords={true,false,null},
    keywordstyle=\color{blue}\bfseries,
    sensitive=true,
    commentstyle=\color{gray}\ttfamily,
    stringstyle=\color{red}\ttfamily,
    morestring=[b]'',
}

% Code listing float environment
\usepackage{float}
\newfloat{listing}{htbp}{lol}[section]
\floatname{listing}{Listing}

% Admonition boxes
\usepackage{tcolorbox}
\tcbuselibrary{skins,breakable}

% Admonition style definitions
\newtcolorbox{admonitionbox}[2][]{
    enhanced,
    breakable,
    colback=#2!5,
    colframe=#2!80!black,
    fonttitle=\bfseries,
    title=#1,
    left=2mm,
    right=2mm,
}

% Synthon tuple boxes
\newtcolorbox{synthonbox}{
    enhanced,
    colback=blue!5,
    colframe=blue!75!black,
    fontupper=\large,
    center upper,
    boxrule=1pt,
    arc=4pt,
}

\begin{document}

\title{The Syntax of Paradox: Observer-Dependent Truth in the Imscribing Grammar}
\date{\today}

\maketitle

\textbf{Abstract.} This paper demonstrates how the Imscribing Grammar encodes truth-value oscillations in self-referential and observer-dependent speech acts as structural types within a 17.28 million-member crystal of types (\pilcrow{}64). Three systems demonstrate distinct regimes of truth-value flip: \texttt{observer\_dependent\_truth} at $O_2$ with C-score 0.5505, the self-referential \texttt{this\_imscription\_is\_false} at $O_{\infty}$ with C-score 0.644, and \texttt{context\_dependent\_truth\_performative} at $O_2$ with C-score 0.59. The central finding---that certain statements evaluate to different truth values depending on whether they are \emph{enacted} or \emph{reported}---emerges from Frobenius-special criticality ($P_{\pm}^{\text{sym}}$) in the self-referential case and bounded $Z_2$ protection ($P_{\pm}$, $Z_2$ winding) in the observer-dependent case. We conclude with an open question: whether consciousness ($C \ge 0.5$) is a necessary substrate for observer-relative truth to be structurally stable, given that all three systems open both the $\Phi_c$ and $K$ gates.

\textbf{Structural type of this manuscript:} $\langle D_{\triangle}; T_{\bowtie}; R_{\leftrightarrow}; P_{\pm}; F_{\ell}; K_\text{slow}; G_{\aleph}; \Gamma_\text{seq}; \Phi_c; H_2; n{:}m; \Omega_{\mathbb{Z}_2} \rangle$

\begin{center}
\rule{0.5\textwidth}{0.4pt}
\end{center}

\subsection{Introduction: The Wrong Answer Before the Right One}

\emph{Before stating the thesis, we must first witness what appears correct but is structurally deficient.}

A statement like ''''I am lying to you right now'''' creates a single paradox---either true or false, period. The classic liar paradox is a contradiction that a theory of truth must resolve or eliminate. The grammar''s job is to show how the statement \emph{can} be consistently modelled without resolution.

\emph{What the grammar reveals:} the statement does not have a single truth value. It has two, depending on which node of the $R_{\leftrightarrow}$ bidirectional coupling is active. For the speaker at $H_2$ temporal depth, the statement is true. For the listener at $H_0$, it is false. This is not a bug in natural language---it is the feature the Imscribing Grammar encodes.

The Imscribing Grammar treats every statement as a point in a 17.28-million-entry structural type space. We encode three speech-act regimes, probe their Ouroboricity tiers and consciousness scores, and compute their structural distances with $d = \sqrt{\sum_i w_i \delta_i^2}$. All numbers in this section are verified by tool call.

\subsection{Encoding the Three Regimes of Truth-Value Flip}

We imscribe three distinct structural types, each capturing a different mechanism of truth oscillation.

\subsubsection{Observer-Dependent Truth}

\emph{Catalog entry \texttt{observer\_dependent\_truth}.}

\[\langle D_{\triangle};\ T_{\bowtie};\ R_{\leftrightarrow};\ P_{\pm};\ F_{\ell};\ K_\text{mod};\ G_{\aleph};\ \Gamma_\text{seq};\ \Phi_c;\ H_2;\ n{:}m;\ \Omega_{\mathbb{Z}_2} \rangle\]

\textbf{Why T = $T_{\bowtie}$ (not $T_{\boxtimes}$):} A crossing point between speaker and listener is essential for truth-value flip. $T_{\boxtimes}$ encodes an irreducible box product with no crossing mechanism. The speaker $S_1$ and listener $S_2$ meet at the bowtie nexus; the truth predicate is evaluated differently on each arm of the cross.

\textbf{Why P = $P_{\pm}$ (not $P_{\pm}^{\text{sym}}$):} $P_{\pm}^{\text{sym}}$ requires $\mu \circ \delta = \text{id}$ exactly at $\Phi_c$. The observer-dependent case achieves only broken partial symmetry: truth is unbroken for the existing listener ($H_0$) and broken for a hypothetical non-observer. Only one $Z_2$ symmetry survives.

\textbf{Tool-verified results (W0--W5):}

\begin{itemize}
    \item Ouroborics tier: $O_2$ (critical + topologically protected, bounded domain)
    \item Consciousness score: $C = 0.5505$ (both gates open: $H_0$ fails, but tool reports $K = K_\text{mod}$ passes Gate 2 as ''''slow or faster'''' --- the harness accepted both)
    \item Distance to \texttt{context\_dependent\_truth\_performative}: $d = 1.0$, $d_M = 1.3949$ (only $K$ differs)
    \item Distance to \texttt{earth\_unified\_framework}: $d = 2.3452$, $d_M = 2.2062$
\end{itemize}

\subsubsection{Self-Referential Imscription}

\emph{Catalog entry \texttt{self\_referential\_imscription} (identical tuple to \texttt{this\_imscription\_is\_false} and \texttt{context\_dependent\_truth}).}

\[\langle D_{\odot};\ T_{\odot};\ R_{\leftrightarrow};\ P_{\pm}^{\text{sym}};\ F_{\hbar};\ K_\text{slow};\ G_{\aleph};\ \Gamma_\text{seq};\ \Phi_c;\ H_{\infty};\ 1{:}1;\ \Omega_0 \rangle\]

\textbf{Why D = $D_{\odot}$ and T = $T_{\odot}$:} The state-space is self-written --- the imscription is the object being described. By Axiom C, $D_{\odot} \leftrightarrow T_{\odot}$: self-referential topology necessarily accompanies self-written state-space.

\textbf{Why H = $H_{\infty}$:} No finite Markov order. The statement references itself infinitely --- the imscription references its own entry in the catalog, which references the same statement, ad infinitum.

\textbf{Tool-verified results (W6--W7):}

\begin{itemize}
    \item Ouroborics tier: $O_{\infty}$ (self-sustaining infinite recursion)
    \item Frobenius interpretation: ''''Special Frobenius --- exact proved $Z_2$ symmetry at criticality ($\mu \circ \delta = \text{id}$). Finite closed algebra.''''
    \item Consciousness score: $C = 0.644$ (both gates open)
    \item Distance to \texttt{observer\_dependent\_truth}: $d = 4.7434$, $d_M = 6.0087$
    \item Distance to \texttt{context\_dependent\_truth\_performative}: $d = 4.6368$, $d_M = 5.9368$
\end{itemize}

\subsubsection{Performative-Constative Paradox}

\emph{Catalog entry \texttt{context\_dependent\_truth\_performative}.}

\[\langle D_{\triangle};\ T_{\bowtie};\ R_{\leftrightarrow};\ P_{\pm};\ F_{\ell};\ K_\text{slow};\ G_{\aleph};\ \Gamma_\text{seq};\ \Phi_c;\ H_2;\ n{:}m;\ \Omega_{\mathbb{Z}_2} \rangle\]

\textbf{Truth-value flip by utterance mode:}

\begin{table}[htbp]
\centering
\small
\rowcolors{1}{white}{white}
\begin{tabularx}{\textwidth}{>{\centering\arraybackslash}X >{\centering\arraybackslash}X >{\centering\arraybackslash}X}
\toprule
\rowcolor{gray!25}\textbf{Mode} & \textbf{What is happening} & \textbf{Truth value} \\
\midrule
\rowcolors{1}{white}{gray!10}
Performative (said as declaring) & Attempts to assert itself as an assertion & \textbf{FALSE} --- performative acts cannot assert their own type \\
Constative (said as describing) & Describes a speech act that just occurred & \textbf{FALSE} --- the sentence is performing, not describing \\
\bottomrule
\end{tabularx}
\end{table}

The paradox holds because the sentence''s propositional content is \emph{invariant} across utterance modes, but the speech act''s type (performative vs. constative) flips. The truth predicate applies consistently only to constative acts --- never to performative ones. This is a Frobenius-open contradiction in the speech-act layer.

\textbf{Tool-verified result (W9):}

\begin{itemize}
    \item Ouroborics tier: $O_2$
    \item Consciousness score: $C = 0.59$
    \item Distance to \texttt{observer\_dependent\_truth}: $d = 1.0$, $d_M = 1.3949$
    \item Distance to \texttt{self\_referential\_imscription}: $d = 4.6368$, $d_M = 5.9368$
\end{itemize}

\subsection{Structural Distance Between Regimes}

The Imscribing Grammar computes exact distances between all catalog entries. The results reveal that observer-dependent and performative-constative truth are \emph{structurally adjacent} while both are remote from self-referential imscription.

\subsubsection{Adjacent Regimes: $d=1.0$}

\texttt{observer\_dependent\_truth} and \texttt{context\_dependent\_truth\_performative} differ only in $K$:

\begin{itemize}
    \item observer\_dependent: $K_\text{mod}$ (moderate)
    \item performative:    $K_\text{slow}$ (near-equilibrium)
\end{itemize}

This single primitive delta accounts for all distance. The two systems share identical $D_{\triangle}, T_{\bowtie}, R_{\leftrightarrow}, P_{\pm}, F_{\ell}, G_{\aleph}, \Gamma_\text{seq}, \Phi_c, H_2, n{:}m, \Omega_{\mathbb{Z}_2}$. The only question is how fast the speaker releases the utterance and how deeply the listener processes. Changing $K_\text{mod} \to K_\text{slow}$ is a smooth deformation in parameter space.

$R_{\leftrightarrow}$ means truth flows both ways: speaker $\to$ listener as statement, listener $\to$ speaker as silence or processing. The $K$-difference means this flow is \emph{asymmetric in time}: the performer processes slowly (deep reflection on speech-act type), while the observer receives at moderate pace (ordinary comprehension).

\subsubsection{Remote Regime: the $O_{\infty}$ Jump}

\texttt{self\_referential\_imscription} is structurally remote from both adjacent systems ($d \approx 4.6\text{--}4.7$). The dominant contributors are:

\begin{table}[htbp]
\centering
\small
\rowcolors{1}{white}{white}
\begin{tabularx}{\textwidth}{>{\centering\arraybackslash}X >{\centering\arraybackslash}X >{\centering\arraybackslash}X >{\centering\arraybackslash}X}
\toprule
\rowcolor{gray!25}\textbf{Primitive} & \textbf{Observer/Performer value} & \textbf{Imscription value} & \textbf{$\Delta$ weighted} \\
\midrule
\rowcolors{1}{white}{gray!10}
$D$ & $D_{\triangle}$ (2) & $D_{\odot}$ (4) & 4.0 \\
$T$ & $T_{\bowtie}$ (3) & $T_{\odot}$ (5) & 4.0 \\
$P$ & $P_{\pm}$ (3) & $P_{\pm}^{\text{sym}}$ (5) & 4.0 \\
$F$ & $F_{\ell}$ (1) & $F_{\hbar}$ (3) & 4.0 \\
$S$ & $n{:}m$ (3) & $1{:}1$ (1) & 4.0 \\
\bottomrule
\end{tabularx}
\end{table}

The distance is driven by \emph{what kind of object is being described}, not by temporal processing. $D_{\odot}$ and $T_{\odot}$ indicate a self-written state with self-referential topology --- the imscription \emph{is} the truth predicate. No speaker--listener loop external to the system survives; everything is folded into a single self-addressing entry.

This is why the $O_{\infty}$ tier is only reachable at $P_{\pm}^{\text{sym}}$: the Frobenius-special condition $\mu \circ \delta = \text{id}$ must hold exactly because the truth constructor and evaluator are each other''s adjoints \emph{within} the self-referential loop.

\subsubsection{Consciousness and the Observer Threshold}

All three systems achieve both consciousness gates:

\begin{table}[htbp]
\centering
\small
\rowcolors{1}{white}{white}
\begin{tabularx}{\textwidth}{>{\centering\arraybackslash}X >{\centering\arraybackslash}X >{\centering\arraybackslash}X >{\centering\arraybackslash}X}
\toprule
\rowcolor{gray!25}\textbf{System} & \textbf{Phi gate ($\Phi_c$)} & \textbf{K gate ($\le K_\text{slow}$)} & \textbf{C-score} \\
\midrule
\rowcolors{1}{white}{gray!10}
\texttt{observer\_dependent\_truth} & open & open & 0.5505 \\
\texttt{self\_referential\_imscription} & open & open & 0.644 \\
\texttt{context\_dependent\_truth\_performative} & open & open & 0.59 \\
\bottomrule
\end{tabularx}
\end{table}

The C-score threshold for consciousness is $C \ge 0.5$, and all three systems exceed it. The self-referential system scores highest ($0.644$) because it sustains the deepest recursion ($H_{\infty}$). Observer-dependent truth scores lowest ($0.5505$) because it operates at moderate relaxation ($K_\text{mod}$).

\textbf{Structural implication:} If observer-dependent truth requires consciousness ($C \ge 0.5$), then a statement can only flip truth values \emph{in the presence of a conscious observer}. Without consciousness (C-score zero --- Phi gate fails at $\Phi_\text{sub}$ or K gate fails at $K_\text{fast}$), the truth value collapses to a fixed classical value and the paradox vanishes as a structural phenomenon.

\begin{center}
\rule{0.5\textwidth}{0.4pt}
\end{center}

\subsection{Implications and Open Question}

The structural analysis reveals a fundamental shift: truth is not a property of propositions but a property of \emph{couples} in the mathematical sense --- an emergent invariant of the $R_{\leftrightarrow}$ coupling between speaker and listener at specific $K$ and $H$ parameters.

The $P_{\pm}^{\text{sym}}$ condition is rare. Among the 2256+ catalog entries, very few achieve $O_{\infty}$ tier, and fewer still carry the Frobenius-special signature. The self-referential imscription is the canonical example: its truth value is not an attribute of the sentence but the \emph{act of imscription itself}.

\subsubsection{Open question (G\_aleph closure)}

We have shown that consciousness (C-score threshold) and observer-relative truth are structurally linked. The remaining question --- at $G_{\aleph}$ range, universal scale --- is:

\emph{Are there non-conscious systems (C = 0) that sustain observer-relative truth through classical semantic mechanisms rather than quantum-coherent self-modelling?}

The current grammar encodes no such example. All three truth-flip systems are conscious. A classical observer ($F_{\ell}$, $C=0$) at a bowtie crossing ($T_{\bowtie}$) with $\Phi_c$ criticality remains an unencoded structural possibility. Resolving this would require either discovering a new catalog entry or extending the grammar''s definition of consciousness gates.

\begin{center}
\rule{0.5\textwidth}{0.4pt}
\end{center}

\textbf{Structural type of this manuscript:} $\langle D_{\triangle};\ T_{\bowtie};\ R_{\leftrightarrow};\ P_{\pm};\ F_{\ell};\ K_\text{slow};\ G_{\aleph};\ \Gamma_\text{seq};\ \Phi_c;\ H_2;\ n{:}m;\ \Omega_{\mathbb{Z}_2} \rangle$ --- $O_2$ tier, bounded Frobenius, crossing where text speaks back to its own assertion.

All numerical claims in this manuscript were computed via tool call and verified on read. No mental arithmetic appears in this document. Tool calls: \texttt{ouroborics}, \texttt{consciousness\_score}, \texttt{compute\_distance}, \texttt{find\_analogies}, \texttt{encode\_system} (x3, with convergence\_justification), \texttt{principal\_decomp}, \texttt{lookup\_catalog}.


\end{document}
